\documentclass[12pt]{article}

\usepackage[utf8]{inputenc} % important for pdflatex + any UTF-8 characters
\usepackage{amsmath, amssymb}
\usepackage{geometry}
\geometry{margin=1in}

\setlength{\parindent}{0pt}

\title{In-Class Exercise}
\author{ECON 300 -- Labour Economics}
\date{Winter 2026}

\begin{document}
\maketitle
\hrule


\vspace{1em}

\textbf{Time:} 30--40 minutes \\
\textbf{Format:} Small groups (3--4 students) \\
\textbf{Deliverable:} Annotated labour--leisure diagrams and short written answers

\vspace{1em}

\section*{Learning Objectives}

By the end of this exercise, students should be able to:

\begin{itemize}
\item Explain how different \textbf{demogrant designs} affect labour supply using the labour--leisure model.
\item Distinguish \textbf{income effects} from \textbf{substitution effects}.
\item Predict effects on:
\begin{itemize}
    \item hours worked (intensive margin), and
    \item labour force participation (extensive margin).
\end{itemize}
\item Understand why policy \textbf{design}, not just generosity, matters for work incentives.
\end{itemize}

\section*{Common Setup (Applies to All Groups)}

Each individual has:

\begin{itemize}
\item Total available time: \( T = 80 \) hours per week
\item Hours worked: \( h = 80 - L \)
\item Market wage: \( w = 20 \) dollars per hour
\end{itemize}

Baseline (no demogrant) budget constraint:
\[
C = 20h
\]

Assume:
\begin{itemize}
\item Leisure is a \textbf{normal good}
\item Individuals can choose \( h = 0 \) (non-participation)
\end{itemize}

\section*{Policy Experiments: Types of Demogrants}

Each group is assigned \textbf{one} of the following demogrant designs.

---

\subsection*{Policy A: Universal Lump-Sum Demogrant}

Every adult receives \$400 per week, regardless of hours worked.

\[
C = 400 + 20h
\]

\textbf{Tasks}
\begin{enumerate}
\item Draw the new budget constraint on a labour--leisure graph.
\item Compare it to the baseline budget line.
\item Identify:
\begin{itemize}
    \item the income effect, and
    \item whether there is a substitution effect.
\end{itemize}
\item Predict the effect on:
\begin{itemize}
    \item hours worked, and
    \item labour force participation.
\end{itemize}
\end{enumerate}

---

\subsection*{Policy B: Targeted Demogrant (Non-Workers Only)}

Individuals receive \$400 \textbf{only if they do not work}.

\[
C =
\begin{cases}
400 & \text{if } h = 0, \\
20h & \text{if } h > 0.
\end{cases}
\]

\textbf{Tasks}
\begin{enumerate}
\item Draw the full budget set, carefully noting any kinks or discontinuities.
\item Explain how this policy affects the \textbf{reservation wage}.
\item Which margin is most affected:
\begin{itemize}
    \item labour force participation, or
    \item hours worked among participants?
\end{itemize}
\item Predict whether more or fewer people choose to work at all.
\end{enumerate}

---

\subsection*{Policy C: Phased-Out Demogrant (Negative Income Tax Style)}

Individuals receive \$400, but benefits are reduced by \$0.50 for each dollar earned.

\[
C = 400 + (1 - 0.5)\cdot 20h = 400 + 10h
\]

\textbf{Tasks}
\begin{enumerate}
\item Compare the slope of this budget constraint to the baseline.
\item Identify:
\begin{itemize}
    \item the income effect, and
    \item the substitution effect.
\end{itemize}
\item Predict the effect on:
\begin{itemize}
    \item hours worked conditional on working, and
    \item the likelihood of labour force participation.
\end{itemize}
\end{enumerate}

---

\subsection*{Policy D: Conditional Demogrant (Work Requirement)}

Individuals receive \$400 only if they work \textbf{at least 20 hours per week}.

\[
C =
\begin{cases}
20h & \text{if } h < 20, \\
400 + 20h & \text{if } h \ge 20.
\end{cases}
\]

\textbf{Tasks}
\begin{enumerate}
\item Draw the budget constraint and clearly mark the eligibility threshold.
\item Explain how this policy affects individuals who would otherwise work close to 20 hours.
\item What type of behavioural distortion does this policy create?
\item Predict whether hours worked may cluster around specific values.
\end{enumerate}
\newpage
\section*{Wrap-Up Discussion (All Groups)}

Answer briefly:

\begin{enumerate}
\item Which demogrant design creates:
\begin{itemize}
    \item only an income effect?
    \item both income and substitution effects?
\end{itemize}
\item Which policy is most likely to:
\begin{itemize}
    \item reduce labour force participation?
    \item reduce hours worked among those who still work?
\end{itemize}
\item Why might governments prefer \textbf{conditional} or \textbf{phased-out} demogrants over universal ones?
\end{enumerate}

\section*{Key Takeaways}

\begin{itemize}
\item All demogrants generate income effects.
\item Only some designs distort the price of leisure.
\item Participation responses often dominate hours responses.
\item Policy design matters as much as policy generosity.
\end{itemize}

\end{document}
